\chapter{Summary and Conclusions}
\label{ch:summary}

\section{Summary of contributions}
\label{sec:sum-contrib}

This thesis has reported a joint-embedding predictive pretraining
recipe for solar magnetic winding-flux cubes at a corpus size of
twenty-one active regions. The principal contributions are
recapitulated below.

A preprocessing pipeline for a signed pseudoscalar field of
variable spatial extent was developed
(Chapter~\ref{ch:data}). The pipeline includes a
chirality-respecting $D_{4}$ augmentation scheme that pairs
spatial flips and quarter-turn rotations with explicit sign
negation; a winding-flux guard of $|W| \le 10^{8}$ established by
the audit of \S\ref{sec:data-clip}; and a cube-level (HARP-identity)
holdout convention that eliminates active-region leakage between
the training, validation, and downstream-probe partitions.

An adaptation of the V-JEPA-2-AC template
(Chapter~\ref{ch:architecture}) to the single-channel
twelve-minute-cadence active-region cube format was specified.
The model uses a ViT-Small context encoder with two-dimensional
rotary positional embeddings in physical (megametre) units, an
exponential-moving-average target encoder, and a six-layer
block-causal predictor with three-dimensional rotary positional
embeddings. A five-family mask catalogue (short tube, long tube,
future block, cross-time, and tail) supplies the masking
distribution; a two-stage curriculum (\S\ref{sec:exp-protocol})
schedules its rollout. The training loss is smooth-$L_{1}$
restricted to the intersection $M_{\text{loss}} = M_{\text{mask}}
\land M_{\text{valid}} \land M_{\text{pad}}$ of the masked,
data-valid, and pad-valid token sets.

Four ablation families fixed the operating point
(Chapter~\ref{ch:experiments}). The mask-policy ablation
(E12--E15) established the uniform mixture as the new sanity-scale
floor at val~$0.00530$, $3.42\times$ better than the tube-only
ablation and $1.57\times$ better than the previous tube-heavy
default. The EMA-decay sweep (E17--E20) established a monotone
monotone response in which a smaller $\tau$ yields lower
validation loss across the tested range, with
$\tau = 0.990$ the winner at val~$0.00761$. The mask-ratio sweep
(E25--E28) established a concave response peaked at $r = 0.75$
(E26, val~$0.00876$). The thesis-scale run (E30) trained on
thirteen cubes for one hundred epochs of two thousand steps each,
achieving a $5.8\times$ throughput improvement over the failed
E29b launch by excluding the five largest cubes from the bucketed
batches that limited throughput.

A frozen-feature wind-flux probe on the E30~v2 thesis-scale
checkpoint (Chapter~\ref{ch:probe}) regressed the spatial-mean
winding-flux curve at $R^{2} = 0.73$ (MLP, $r = 0.87$, medAPE
$7.3\%$) on the encoder-validation cubes and remained positive on
five novel cubes that the encoder never saw ($R^{2} = +0.23$
aggregate, $r = +0.50$). A per-cube log-affine recalibration with
one $(a, b)$ pair recovered a median medAPE of $9.9\%$ across eight
evaluation cubes and beat the lag-one persistence baseline on
four of five novel cubes (Finding~F9, confirmed). This is the
primary downstream evidence of the present work that the JEPA
encoder transfers as a wind-flux feature extractor to unseen
active regions.

A frozen-feature binary flare classifier was attached to the same
cached features under a within-cube temporal split
(\S\ref{sec:probe-flare}). The M$+$/24~h configuration was bounded
by the lag-one persistence floor at the present corpus size: only
one informative evaluation cube exists at this class--window pair,
and the head AUC of $0.252$ trails the persistence TSS of $0.973$.
A C$+$/12~h two-layer MLP head delivered AUC $= 1.000$ and
TSS $= 1.000$ on harp\_49, exceeding the persistence floor of
$0.975$ on the densest mixed-label evaluation cube and constituting
the first frozen-feature flare classifier in the present work to
beat persistence on an informative cube (Finding~F13). The result
is conditional rather than uniform: of eight mixed-label cubes
under the C$+$ class, the harp\_49 positive is the only configuration
to beat persistence; no head--window pair dominates the cross-cube
aggregate at twenty-one cubes.

\section{Limitations}
\label{sec:sum-limits}

The principal limitations of the present work are listed below in
decreasing order of severity.

\paragraph{Corpus scale.} Twenty-one active-region cubes is two
orders of magnitude smaller than the smallest validated video
self-supervised regime in the
literature~\cite{tong2022videomae,wang2023videomaev2}. The
sanity-floor losses reported in Chapter~\ref{ch:experiments}
characterise the optimisation behaviour of the present recipe but
do not establish that the resulting encoder transfers to a
downstream task at the level of larger-corpus foundation models.

\paragraph{Pixel-only ceiling.} Even with a perfect downstream
head, pixel-only flare forecasters are bounded by a true-skill-
statistic ceiling of approximately~$0.7$. The fusion of pixel
features with photospheric scalar summaries (SHARP
parameters~\cite{bobra2014sharp}) is required to break past this
ceiling and is deferred to a successor configuration.

\paragraph{Downstream-head scope.} The wind-flux probe of
\S\S\ref{sec:probe-stage1}--\ref{sec:probe-calibration} is a
regression diagnostic anchored to lag-one persistence, not an
operational forecast. The binary flare classifier of
\S\ref{sec:probe-flare} is anchored to lag-one persistence and
reports the true skill statistic at the per-cube level; it
attains a positive result on harp\_49 at the C$+$/12~h pair and
does not generalise uniformly across the eight mixed-label
evaluation cubes. Critical success index and Heidke skill score
on a pixel-space reconstructed field require a decoder head and
a GOES-event catalogue merge, neither of which is implemented in
the present work.

\paragraph{Single backend coupling.} The Apple~MPS backend
requires a hand-written attention path under
\texttt{torch.no\_grad} because the SDPA kernel returns NaN with
non-trivial \texttt{attn\_mask} (Finding~F6). The CUDA and MPS
paths are not unified; unifying them would silently regress MPS
evaluation. The coupling is benign but increases the maintenance
surface.

\paragraph{Patch-size and temporal-context unswept.} The patch
size, $t_{\text{in}}$, $t_{\text{out}}$, and predictor depth were
fixed at the values reported in Chapter~\ref{ch:architecture} and
not swept; downstream sensitivity to these settings is unknown.

\section{Future outlook}
\label{sec:sum-future}

The roadmap from the present pretraining result to operational
flare forecasting comprises four steps.

\paragraph{Strategy A: visible-only encoder.} The current
implementation processes all tokens, including masked-out ones,
through the context encoder. A true V-JEPA visible-only encoder
processes only the visible tokens, recovering the
$1{-}r \approx 25\%$ computational fraction and enabling a
larger predictor for the same wall-clock budget. The rewrite
requires a sparse-token patchifier and a corresponding predictor
contract change.

\paragraph{Pixel decoder.} A
DPT~\cite{ranftl2021dpt}- or
SegFormer~\cite{xie2021segformer}-style decoder attached to the
frozen encoder would produce pixel-space reconstructions of the
winding-flux field and supply persistence-anchored Critical
Success Index (CSI) and Heidke Skill Score (HSS) numbers. The decoupling principle of Chapter~\ref{ch:probe}
applies: the decoder composes freely with the frozen encoder and
does not require pretraining changes.

\paragraph{Pretraining at scale on SuryaBench.} The SuryaBench
corpus~\cite{suryabench2025} is the natural next pretraining
substrate. The SuryaBench data are at $4096^{2}$ spatial resolution
with thirteen channels at sixty-minute cadence, so the present
single-channel twelve-minute pipeline would require a channel
adapter and a temporal interpolation; both are well-posed
modifications that do not touch the joint-embedding doctrine.

\paragraph{Multimodal SHARP-GOES fusion (V5.2).} Pixel features
fused with the SHARP photospheric scalar
summaries~\cite{bobra2014sharp} have been reported in the
published literature to exceed the pixel-only TSS ceiling,
breaking past the limit observed for pixel-only inputs. A separate pipeline merging cube
features with SHARP scalar timeseries and GOES-event labels is the
target configuration for a successor work; it is outside the scope
of the present thesis.

\section{Concluding remark}
\label{sec:sum-final}

The joint-embedding predictive doctrine adapts to the solar
magnetic-winding-flux setting at a corpus size two orders of
magnitude below the smallest validated video self-supervised
regime cited in this thesis. The wind-flux probe of
Chapter~\ref{ch:probe} is the load-bearing evidence: the frozen
E30~v2 encoder regresses $\langle W \rangle(t)$ on five novel
active regions to median absolute percentage error $9.9\%$ after
a two-parameter per-cube recalibration, beating lag-one
persistence on four of them (Finding~F9). A secondary
frozen-feature flare classifier reaches AUC~$=1.000$ on the
densest mixed-label evaluation cube under the C$+$/12~h
configuration (Finding~F13), exceeding the persistence floor on
the only cube that admits the test at the present corpus size.
The thesis hypothesis --- that JEPA pretraining is feasible at
small-corpus solar scale and that the resulting features carry
the literature-supported flare-precursor signal --- is, within
the disclosed calibration caveat, supported. The
cross-active-region generalisation that would constitute an
operational forecaster remains future work; the next
configuration is a larger pretraining substrate and a multimodal
downstream head, and the present configuration is the foundation
on which that successor will rest.
